\section{Divide and Conquer DP Optimization}
\whenuse{DP por particiones $dp[k][i]=\min_{j<i}(dp[k-1][j]+cost(j,i))$ cuando
el punto \`optimo es mon\'otono en $i$.}

Para transiciones
\[
dp_{new}[m]=\min_{0\le k\le m}\{dp_{old}[k]+cost(k,m)\},
\]
cuando el \'indice \'optimo es mon\'otono: $opt[m]\le opt[m+1]$.

\begin{lstlisting}[style=cpp]
template<class Cost>
void compute_divide_conquer(int left, int right,
                            int option_left, int option_right,
                            const vector<ll> &previous,
                            vector<ll> &current, Cost cost) {
    if (left > right) return;
    while (option_left <= option_right && previous[option_left] == LINF)
        ++option_left;
    while (option_left <= option_right && previous[option_right] == LINF)
        --option_right;
    if (option_left > option_right || right < option_left) {
        fill(current.begin() + left, current.begin() + right + 1, LINF);
        return;
    }
    int middle = left + (right - left) / 2;
    pair<ll,int> best{LINF, -1};
    int last_option = min(middle, option_right);
    for (int option = option_left; option <= last_option; ++option) {
        if (previous[option] == LINF) continue;
        ll transition = cost(option, middle);
        if (transition == LINF) continue;
        i128 raw = (i128)previous[option] + transition;
        ll candidate = raw > LINF ? LINF : raw < -LINF ? -LINF : (ll)raw;
        best = min(best, pair<ll,int>{candidate, option});
    }
    current[middle] = best.first;
    // Sin optimo finito no hay informacion para podar a los vecinos.
    int left_end = best.second == -1 ? option_right : best.second;
    int right_begin = best.second == -1 ? option_left : best.second;
    compute_divide_conquer(left, middle - 1, option_left, left_end,
                           previous, current, cost);
    compute_divide_conquer(middle + 1, right, right_begin, option_right,
                           previous, current, cost);
}
\end{lstlisting}

La monoton\'ia debe demostrarse (por ejemplo con desigualdad cuadrangular o
matriz Monge); no alcanza con que el costo ``parezca convexo''.
\texttt{LINF} representa estados o transiciones imposibles; no se suma a un
costo finito. Si no hay candidato finito, el estado queda en \texttt{LINF} y
ambos lados conservan los l\'imites de candidatos. Los costos y sumas finitos
deben estar estrictamente entre $-\texttt{LINF}$ y $\texttt{LINF}$.

\complexity{$O(kn\log n)$ para $k$ capas si evaluar \texttt{cost} cuesta
$O(1)$ y cada estado procesado tiene candidato finito; memoria $O(n)$ con dos
capas. Si hay estados imposibles por restricciones de \texttt{cost}, el
fallback sin poda puede costar $O(kn^2)$; conviene acotar el dominio a los
estados factibles cuando el problema lo permita.}

\subsection{Acotar los estados factibles}

Para dividir los primeros $i$ elementos en exactamente $g$ segmentos no
vacíos, los estados $i<g$ son imposibles. Si el último segmento es $[j,i)$,
el corte cumple $g-1\le j<i$. Inicializar los demás estados en
\texttt{LINF} y calcular sólo $i\in[g,n]$:

\begin{lstlisting}[style=cpp]
template<class SegmentCost>
ll partition_nonempty(int n, int groups, SegmentCost segment_cost) {
    assert(n >= 0);
    if (groups < 0 || groups > n) return LINF;
    vector<ll> previous(n + 1, LINF), current(n + 1, LINF);
    previous[0] = 0; // Cero elementos y cero grupos.
    for (int g = 1; g <= groups; ++g) {
        fill(current.begin(), current.end(), LINF);
        auto cost = [&](int j, int i) -> ll {
            if (j >= i) return LINF; // Segmento vacio o invertido.
            return segment_cost(j, i); // Costo de [j,i).
        };
        compute_divide_conquer(g, n, g - 1, n - 1,
                               previous, current, cost);
        previous.swap(current);
    }
    return previous[n];
}
\end{lstlisting}

Si todo segmento no vacío tiene costo finito, cada estado procesado tiene al
menos un candidato finito: el prefijo de longitud $g-1$ puede formar $g-1$
grupos y dejar el resto como último segmento. Con óptimos monótonos y costo
$O(1)$, se conserva $O(\texttt{groups}\cdot n\log n)$ tiempo y $O(n)$ memoria.
No hace falta que todas las transiciones sean válidas.

\warning{Que existan exactamente $g$ grupos no garantiza la monotonía del
corte óptimo; hay que demostrarla para el costo del problema. Si hay otras
restricciones sobre los segmentos, pueden quedar estados imposibles dentro
de $[g,n]$ y esta llamada por sí sola no elimina el fallback cuadrático.}

Si se conocen estados factibles con huecos, se puede adaptar la recursión
para recorrer una lista ordenada de ellos. Conservar los índices originales
al evaluar costos y acotar candidatos; la posición dentro de la lista no es
el extremo del segmento. Si la factibilidad sólo se descubre probando todos
los candidatos, esta plantilla no garantiza evitar $O(n^2)$ por capa.
